\subsection{Fazit}
\label{subsec:fazit}

Die in Abschnitt~\ref{subsec:zielsetzung} formulierte Zielsetzung wurde
hinsichtlich der strukturierten technischen Untersuchung und qualitativen
Bewertung weitgehend erreicht. Die Fallstudie zeigt, dass ein gemeinsamer
Kern, deklarative Profile, getrennte Capabilities und ein zentraler
Tooling-Einstieg eine konsistente, wiederverwendbare Baseline für die
definierten Web-, Cloud- und Desktopvarianten bilden können.

Der wesentliche Nutzen liegt in kontrollierter Variabilität und der
Vereinheitlichung wiederkehrender Projekt- und Entwicklungsabläufe. Daraus
ergibt sich ein plausibles Produktivitätspotenzial, dessen tatsächliche Größe
ohne empirische Zeit- und Aufwandsdaten offen bleibt. Ebenso ist die
Übertragung auf Kundenprojekte bislang nur für den technisch abgegrenzten
Zielbereich analytisch begründbar.

Das Gesamturteil endet daher bei der Reife einer untersuchten technischen
Template-Baseline. Die lokale Abnahme stützt den dafür betrachteten Stand;
aktuelle CI-Fehler, offene externe Prüfungen und produktspezifische
Auslieferungs- und Sicherheitsaufgaben verhindern jedoch eine Gleichsetzung
mit einem fertigen, plattformübergreifend verifizierten Produkt. Gerade diese
Trennung bestimmt die Aussagekraft der Fallstudie.
